\section{The Fifth Path}

Today, all young people are witness to and should be aware of the rising tide of `burnt out,' thoroughly dysfunctional, and often mentally ill individuals who surround them. This is most obvious in university environments, where arbitrary expressions of `individualism' are highly encouraged and normalcy is ruthlessly attacked as old-fashioned, unnatural, or otherwise boring. A new `normalcy' has been created in these environments that is devoid of any sane point of reference, and since the great majority of people desire to conform to what they believe is normal, a social landscape that is especially difficult to navigate is created.

In fragmented societies, the sense of self is thrust into a very precarious position. Generally, people wish to be noticed and acknowledged by their peers. When recognition does not take place, a questioning of the self takes place and one either reforms one's character or seeks out alternative social circles. In today's world though, a world that has reached a height of fragmentation, participation in society has come to require an excessive pruning one's social image and a need to project an image of power in the grossest material sense of the word. Cultivation of one's faculties or relationships is highly discouraged and looked upon suspiciously by others. Instead, one is expected to have a minor interest in everyone and everything.

Another highly significant factor contributing to the social disorder is the tendency for individuals to live on the basis of their shallowest emotions. For many, this trait has been thoroughly taught since the earliest days of schooling, and more prudent ways of living are labeled `cold,' and `inhuman.' A general desire for people to act rationally or at least responsibly is nothing short of fascist.

Maintaining one's integrity and sense of self in this environment makes you a criminal in the eyes of many. There is extreme form of resentment towards independent types that masks itself with claims to respecting individuality. If you have a sense of irony, you will find it humorous that the preachers of tolerance generally will try to reform everything about your character if you let them.

To return to the more general topic of society, four classes have arisen in modern society, and one can of course be a hybrid of any of these four:

\begin{enumerate}
\item Those who `play the game' and move from place to place flaunting their ``power'' along the way. Their aim is to acquire as much social capital as quickly as possible. The pace here is set to such a high speed that any time for reflection of what they are doing is impossible, and this comes often with great expense to mental and physical well-being.

\item Those who distance themselves from this and try to maintain a sense of self and normalcy as much as is permissible. This is more difficult than it seems as these types are ruthlessly discriminated against by those in positions of power and are increasingly being prepared for and often accepting a slave-like existence in today's social system, simultaneously carrying society on their backs yet deriving an increasingly small pleasure from it. Unfortunately, one is no longer safe here, as some of the more extreme perversions of the day (mixing across castes, pornography, certain types of drugs) have become mainstream.

\item Those who whether from a lack of interior vitality, sheer confusion about how they should act, overwhelming, conscious or unconscious depression about the state of life in general, or actual perversion `join the cesspool' and accept a life of the basest material indulgence without any regard for their sense of self, inner integrity, and their future or that of the world at large. It is astonishing and quite comical if one has a sense of humor about such things that this sort of living is often projected as having moral worth, even to the point of being attributed the highest kind of virtue.

\item Those who refuse to accept it all any longer and distance themselves or remove themselves from society altogether. These types are extremely unpredictable. They can develop in all sorts of ways. In more unfortunate manifestations, they can become a danger to themselves or others because they are unintelligent and can't deal with such an extremely stressful situation, extreme misfortune develop in their life, or they are acutely aware of how things are and willingly become psychopaths. On the flipside, there is little hope for improvement in the general situation except from these types of people, and for those who have a vested interest in things continuing this way, this makes them all the more dangerous.
\end{enumerate}

Clearly, this has become a difficult situation to navigate for those who have not shut their minds off and can look at things lucidly. As always, Gornahoor suggests readers to turn to traditional knowledge as the source of how to live their lives, as it literally is the source of life, healthy civilization, and all that the well-bred man knows to be good. In the Catholicism, there are still some who have at least an inkling of sense of the general disarray and a will and rites to create healthier human ways of life, but one will still have to dig deeper into the Tradition to rise above this situation, as much in the exoteric Church has been lost to the profane world.


\flrightit{Posted on 2011-02-05 by VisionsOfGlory14}

\begin{center}* * *\end{center}

\begin{footnotesize}
\begin{sffamily}
\texttt{Lovekraft on 2011-02-06 at 16:20 said:}

I just found this article from a link in InMaleFide.

Academia as a bastion of free inquiry is a myth. Political correctness, white guilt, and anti-Westernism are the flavors of the day and anyone right-of-center is under constant attack.

But, the Mens Rights Movement is growing into a force to be reckoned with because it has attracted the intelligent disenfranchised. We don't care about the superficial and the greedy and we know that in time, they will come to us for guidance.

\hfill

\texttt{Cologero on 2011-02-11 at 18:12 said:}

In ``Spiritual Authority and Temporal Power'', Rene Guenon alludes to some bad intellectual habits ``acquired through an education that today more often than not produces real mental deformity.''

\hfill

\texttt{kadambari on 2011-02-11 at 22:23 said:}

``Political correctness, white guilt, and anti-Westernism are the flavors of the day and anyone right-of-center is under constant attack.''

I do not think it ``white guilt'' or ``anti-Western'' as people claim although you are right it is politically correct, but the liberal education is controlled by the ivy league liberal elite and they pretty much produce the same generic kind of thinking that is often easy to predict and shallow. If you graduate from such institutions, whatever you pen is likely to be published from what I have seen.

Education has become a thoroughly mercantile venture funded and controlled by the mercantile castes who promote their mercantile values, and there is little in the way of real ``free'' thought in reality.

\hfill

\texttt{Perennial on 2011-02-12 at 02:22 said:}

Excellent reflections, Kadambari. You obviously have had more than a little exposure to the ``educated'' classes.

\hfill

\texttt{Jupiter on 2011-02-26 at 11:49 said:}

``On the flipside, there is little hope for improvement in the general situation except from these types of people, and for those who have a vested interest in things continuing this way, this makes them all the more dangerous.''

I believe this statement indicates that further elaboration on the fourth category is necessary at some point in the future.


\hfill

\end{sffamily}
\end{footnotesize}

